% !TeX encoding = UTF-8
% !TeX TS-program = xelatex
% 车辆运动方程课堂汇报：独立正文稿
% 基于仓库现有 SWJTU Beamer 主题；不修改母版文件。
% 仅包含正文区域：无封面、页眉、导航栏、页面标题、页脚及页码。
% Compile from repository root:
%   xelatex -interaction=nonstopmode -halt-on-error vehicle-motion-equations.tex
%   xelatex -interaction=nonstopmode -halt-on-error vehicle-motion-equations.tex

\documentclass[UTF8,aspectratio=169]{beamer}

\usetheme[
  ColorDisplay=BSblue,
  Miniframes=none,
  ContentMuticols=false,
  BIBMode=none,
  Background=false,
  ResearchCardTitleAlign=left,
  SectionTOC=hide
]{swjtu}

\usepackage{amsmath,amssymb,mathtools}
\usepackage{tikz}
\usepackage[most]{tcolorbox}
\usetikzlibrary{arrows.meta,calc,positioning,decorations.pathmorphing,patterns}

% Keep the existing SWJTU theme, but suppress all outer-matter elements
% for these body-only slides.
\setbeamertemplate{headline}{}
\setbeamertemplate{footline}{}
\setbeamertemplate{frametitle}{}
\setbeamertemplate{navigation symbols}{}
\setbeamertemplate{background}{}
\setbeamercolor{background canvas}{bg=white}
\setbeamersize{text margin left=.7cm,text margin right=.7cm}

% Main body typography is 22 pt or larger. Citation text is intentionally
% secondary but remains comfortably readable on the full-size PPT canvas.
\swjtusetfonts{
  body=22pt,body-leading=27pt,
  card-title=25pt,card-title-leading=29pt,
  diagram=22pt,chart-label=20pt,chart-title=22pt,
  caption=18pt,table=21pt,
  large=27pt,larger=31pt,display=35pt
}
\renewcommand{\normalsize}{\swjtufont{body}}
\colorlet{PanelFill}{AuxiliaryC!5!white}
\colorlet{PanelHead}{AuxiliaryC!10!white}
\colorlet{PanelBorder}{AuxiliaryC!55!white}
\colorlet{SoftOrange}{IntersperseC!10!white}
\colorlet{SoftGreen}{SecondaryAuxiliaryC!8!white}

\newtcolorbox{PPTPanel}[1]{
  enhanced,
  colback=PanelFill,
  colframe=PanelBorder,
  colbacktitle=PanelHead,
  coltitle=PrimaryC,
  title={#1},
  fonttitle=\swjtufont{card-title}\bfseries,
  boxrule=.65pt,
  arc=1.6mm,
  left=2.5mm,right=2.5mm,
  top=1.4mm,bottom=1.6mm,
  before skip=0pt,after skip=0pt
}
\newenvironment{BottomBand}{
  \begin{tcolorbox}[
    enhanced,
    colback=AuxiliaryC!7!white,
    colframe=AuxiliaryC!45!white,
    boxrule=.65pt,arc=1.5mm,
    left=3mm,right=3mm,top=2mm,bottom=2mm
  ]
}{
  \end{tcolorbox}
}
\newenvironment{SlideBody}{
  \begin{frame}[plain,noframenumbering]
  \vspace*{1.8cm}
  \begin{minipage}[t][13.65cm][t]{\linewidth}
  \vspace{0pt}
}{
  \end{minipage}
  \end{frame}
}
\newcommand{\PanelBulletsStart}{%
  \begin{itemize}
  \setlength{\itemsep}{.12cm}
  \setlength{\topsep}{.08cm}
  \setlength{\parskip}{0pt}
}
\newcommand{\PanelBulletsEnd}{\end{itemize}}
\newcommand{\BlueArrow}{\textcolor{AuxiliaryC}{\(\boldsymbol{\longrightarrow}\)}}
\newcommand{\ForceColor}[1]{\textcolor{IntersperseC}{#1}}
\newcommand{\KeyColor}[1]{\textcolor{SecondaryAuxiliaryC}{#1}}

\begin{document}

% ============================================================
% Section 1: 基本知识体系
% ============================================================
\section{基本知识体系}
\subsection{车辆结构与运动自由度}

% P1 — 车辆结构体系与动力传递路径
\begin{SlideBody}
\begin{columns}[T,onlytextwidth]
  \begin{column}{.37\textwidth}
    \begin{minipage}[t][10.65cm][s]{\linewidth}
      \begin{PPTPanel}{车辆动力学组成}
        \PanelBulletsStart
          \item 车体：承载质量与车体振动
          \item 转向架：支承、导向与载荷分配
          \item 轮对：轮轨接触与动力输入
        \PanelBulletsEnd
      \end{PPTPanel}
      \vfill
      \begin{PPTPanel}{两级悬挂系统}
        \PanelBulletsStart
          \item 一系悬挂：轮对 \(\leftrightarrow\) 转向架
          \item 二系悬挂：转向架 \(\leftrightarrow\) 车体
        \PanelBulletsEnd
      \end{PPTPanel}
      \vfill
      \begin{PPTPanel}{主要动力学参数}
        质量与转动惯量；悬挂刚度 \(k\)；阻尼 \(c\)。
      \end{PPTPanel}
    \end{minipage}
  \end{column}
  \begin{column}{.61\textwidth}
    \begin{PPTPanel}{车辆结构与悬挂示意}
      \centering
      \begin{tikzpicture}[x=.98cm,y=.98cm,>=Stealth,line cap=round,line join=round]
        % Car body
        \fill[PanelFill] (.5,5.7) rounded rectangle (18.0,7.0);
        \draw[AuxiliaryC,line width=1.2pt] (.5,5.7) rounded rectangle (18.0,7.0);
        \foreach \x in {1.0,2.8,4.6,6.4,8.2,10.0,11.8,13.6,15.4,17.0}{
          \filldraw[fill=white,draw=AuxiliaryC,line width=.7pt]
            (\x,6.15) rectangle ++(1.15,.48);
        }
        \node[font=\swjtufont{diagram}\bfseries,text=PrimaryC] at (9.2,6.65) {车体};
        % Secondary suspensions
        \foreach \x in {3.7,5.0,13.3,14.6}{
          \draw[IntersperseC,line width=1.2pt]
            (\x,5.68)--(\x,5.47)--++(.16,-.15)--++(-.32,-.18)
            --++(.32,-.18)--++(-.32,-.18)--++(.16,-.15)--(\x,4.43);
        }
        \draw[MutedTextC,line width=1pt] (2.5,4.35) rectangle (6.2,4.65);
        \draw[MutedTextC,line width=1pt] (12.1,4.35) rectangle (15.8,4.65);
        \node[font=\swjtufont{diagram},text=NormalTextC] at (9.2,4.0) {转向架};
        % Primary suspensions and wheels
        \foreach \x in {3.0,4.4,12.6,14.0}{
          \draw[SecondaryAuxiliaryC,line width=1.2pt]
            (\x,4.33)--(\x,4.12)--++(.14,-.14)--++(-.28,-.16)
            --++(.28,-.16)--++(-.28,-.16)--++(.14,-.14)--(\x,3.2);
        }
        \foreach \x in {2.8,4.5,12.4,14.1}{
          \filldraw[fill=white,draw=PrimaryC,line width=1.1pt] (\x,2.85) circle (.42);
          \filldraw[fill=PanelFill,draw=AuxiliaryC,line width=.7pt] (\x,2.85) circle (.16);
        }
        \draw[NormalTextC,line width=2pt] (.5,2.25)--(18,2.25);
        \foreach \x in {1.0,2.2,...,17.6}{
          \draw[MutedTextC,line width=.9pt] (\x,2.1)--++(.38,.3);
        }
        \node[font=\swjtufont{diagram},text=SecondaryAuxiliaryC] at (9.2,1.55) {轮轨接触};
        \draw[->,IntersperseC,line width=1.2pt] (9.2,1.72)--(9.2,2.55);
        \node[font=\swjtufont{caption},text=IntersperseC,anchor=west] at (9.5,2.0) {轮轨力};
        \node[font=\swjtufont{caption},text=NormalTextC] at (2.2,5.0) {二系悬挂};
        \node[font=\swjtufont{caption},text=NormalTextC] at (2.1,3.65) {一系悬挂};
        \node[font=\swjtufont{caption},text=NormalTextC] at (16.4,3.0) {轮对};
      \end{tikzpicture}
    \end{PPTPanel}
  \end{column}
\end{columns}
\vfill
\begin{BottomBand}
  \centering
  \textcolor{PrimaryC}{\bfseries 动力传递路径}\quad
  轮轨力 \(\leftrightarrow\) 轮对 \(\leftrightarrow\) 一系悬挂
  \(\leftrightarrow\) 转向架 \(\leftrightarrow\) 二系悬挂
  \(\leftrightarrow\) 车体
\end{BottomBand}
\end{SlideBody}

% P2 — 车辆运动自由度
\subsection{车辆自由度与建模方法}

\begin{SlideBody}
\begin{columns}[T,onlytextwidth]
  \begin{column}{.36\textwidth}
    \begin{minipage}[t][10.55cm][s]{\linewidth}
      \begin{PPTPanel}{平移自由度}
        \PanelBulletsStart
          \item 纵向 \(x\)：沿线路方向
          \item 横向 \(y\)：垂直线路方向
          \item 竖向 \(z\)：沉浮运动
        \PanelBulletsEnd
      \end{PPTPanel}
      \vfill
      \begin{PPTPanel}{转动自由度}
        \PanelBulletsStart
          \item 侧滚 \(\phi\)：绕纵轴
          \item 点头 \(\theta\)：绕横轴
          \item 摇头 \(\psi\)：绕竖轴
        \PanelBulletsEnd
      \end{PPTPanel}
      \vfill
      \begin{PPTPanel}{自由度选择}
        由研究方向、目标响应和计算成本共同决定。
      \end{PPTPanel}
    \end{minipage}
  \end{column}
  \begin{column}{.62\textwidth}
    \begin{PPTPanel}{车辆空间运动的自由度}
      \centering
      \begin{tikzpicture}[x=1cm,y=1cm,>=Stealth,line cap=round,line join=round]
        % 3D body/cube
        \fill[PanelFill] (6.2,3.4)--(9.0,3.4)--(10.0,4.2)--(7.2,4.2)--cycle;
        \fill[AuxiliaryC!12] (6.2,3.4)--(6.2,5.2)--(7.2,6.0)--(7.2,4.2)--cycle;
        \fill[AuxiliaryC!18] (6.2,5.2)--(9.0,5.2)--(10.0,6.0)--(7.2,6.0)--cycle;
        \draw[PrimaryC,line width=1.3pt] (6.2,3.4)--(9.0,3.4)--(10.0,4.2)--(10.0,6.0)--(7.2,6.0)--(6.2,5.2)--cycle;
        \draw[PrimaryC,line width=1.1pt] (6.2,5.2)--(9.0,5.2)--(10.0,6.0);
        \draw[PrimaryC,line width=1.1pt] (9.0,3.4)--(9.0,5.2);
        \node[font=\swjtufont{diagram}\bfseries,text=PrimaryC] at (8.0,4.65) {车辆刚体};
        % translation axes
        \draw[->,AuxiliaryC,line width=1.5pt] (6.2,3.4)--(3.5,3.4)
          node[left,font=\swjtufont{diagram},text=AuxiliaryC] {\(x\)};
        \draw[->,AuxiliaryC,line width=1.5pt] (6.2,3.4)--(7.7,2.0)
          node[right,font=\swjtufont{diagram},text=AuxiliaryC] {\(y\)};
        \draw[->,AuxiliaryC,line width=1.5pt] (6.2,3.4)--(6.2,1.0)
          node[below,font=\swjtufont{diagram},text=AuxiliaryC] {\(z\)};
        % rotational arrows
        \draw[->,IntersperseC,line width=1.3pt] (4.0,5.5) arc[start angle=210,end angle=500,radius=.65];
        \node[font=\swjtufont{diagram},text=IntersperseC] at (3.8,6.25) {\(\phi\)};
        \draw[->,IntersperseC,line width=1.3pt] (11.3,4.8) arc[start angle=40,end angle=330,radius=.62];
        \node[font=\swjtufont{diagram},text=IntersperseC] at (12.1,5.45) {\(\theta\)};
        \draw[->,IntersperseC,line width=1.3pt] (11.3,2.3) arc[start angle=25,end angle=315,radius=.62];
        \node[font=\swjtufont{diagram},text=IntersperseC] at (12.1,1.9) {\(\psi\)};
        \node[font=\swjtufont{body},text=NormalTextC,align=center] at (8.0,.25)
          {三个平移自由度 \quad 三个转动自由度};
      \end{tikzpicture}
    \end{PPTPanel}
  \end{column}
\end{columns}
\vfill
\begin{BottomBand}
  \centering
  自由度选择 \(\longrightarrow\) 可描述的车辆运动 \(\longrightarrow\) 车桥耦合响应的预测能力
\end{BottomBand}
\end{SlideBody}

% P3 — 车辆动力学模型分类与比较
\subsection{车辆动力学建模与方程建立}

\begin{SlideBody}
\begin{columns}[T,onlytextwidth]
  \begin{column}{.245\textwidth}
    \begin{PPTPanel}{移动荷载}
      \centering
      \begin{tikzpicture}[x=.82cm,y=.82cm,>=Stealth]
        \draw[NormalTextC,line width=1.5pt] (.2,1.1)--(7.0,1.1);
        \draw[NormalTextC,line width=1.5pt] (.5,.65)--(.9,1.1)--(1.3,.65);
        \draw[NormalTextC,line width=1.5pt] (5.9,.65)--(6.3,1.1)--(6.7,.65);
        \draw[->,IntersperseC,line width=1.5pt] (2.5,3.3)--(2.5,1.35);
        \draw[->,IntersperseC,line width=1.5pt] (4.5,3.3)--(4.5,1.35);
        \node[font=\swjtufont{diagram},text=IntersperseC] at (3.5,3.65) {\(P(t)\)};
      \end{tikzpicture}
      \PanelBulletsStart
        \item 给定轴重
        \item 无车辆振动自由度
      \PanelBulletsEnd
      \vfill
      {\swjtufont{body}\bfseries 适用：}\par
      初步桥梁动力计算
    \end{PPTPanel}
  \end{column}
  \begin{column}{.245\textwidth}
    \begin{PPTPanel}{移动质量}
      \centering
      \begin{tikzpicture}[x=.82cm,y=.82cm,>=Stealth]
        \fill[PanelFill] (2.2,2.0) rectangle (4.5,3.5);
        \draw[AuxiliaryC,line width=1.2pt] (2.2,2.0) rectangle (4.5,3.5);
        \node[font=\swjtufont{diagram},text=PrimaryC] at (3.35,2.75) {\(m\)};
        \draw[->,IntersperseC,line width=1.4pt] (3.35,1.9)--(3.35,1.2);
        \draw[NormalTextC,line width=1.5pt] (.2,1.1)--(7.0,1.1);
        \draw[NormalTextC,line width=1.5pt] (.5,.65)--(.9,1.1)--(1.3,.65);
        \draw[NormalTextC,line width=1.5pt] (5.9,.65)--(6.3,1.1)--(6.7,.65);
      \end{tikzpicture}
      \PanelBulletsStart
        \item 考虑移动惯性
        \item 无悬挂系统描述
      \PanelBulletsEnd
      \vfill
      {\swjtufont{body}\bfseries 适用：}\par
      移动质量效应分析
    \end{PPTPanel}
  \end{column}
  \begin{column}{.245\textwidth}
    \begin{PPTPanel}{简化车辆模型}
      \centering
      \begin{tikzpicture}[x=.82cm,y=.82cm,>=Stealth]
        \fill[PanelFill] (2.0,3.1) rectangle (4.8,4.0);
        \draw[AuxiliaryC,line width=1.1pt] (2.0,3.1) rectangle (4.8,4.0);
        \node[font=\swjtufont{diagram}] at (3.4,3.55) {\(m_s\)};
        \draw[IntersperseC,line width=1.2pt] (2.7,3.1)--(2.7,2.85)--++(.18,-.12)--++(-.36,-.16)--++(.36,-.16)--++(-.36,-.16)--(2.7,2.15);
        \draw[SecondaryAuxiliaryC,line width=1.2pt] (4.0,3.1)--(4.0,2.8)--(3.8,2.8)--(3.8,2.35)--(4.2,2.35)--(4.2,2.8)--(4.0,2.8)--(4.0,2.15);
        \fill[PanelFill] (2.0,1.25) rectangle (4.8,2.05);
        \draw[AuxiliaryC,line width=1.1pt] (2.0,1.25) rectangle (4.8,2.05);
        \node[font=\swjtufont{diagram}] at (3.4,1.65) {\(m_u\)};
        \draw[NormalTextC,line width=1.5pt] (.2,.65)--(7.0,.65);
      \end{tikzpicture}
      \PanelBulletsStart
        \item 保留悬挂动力学
        \item 少自由度、效率高
      \PanelBulletsEnd
      \vfill
      {\swjtufont{body}\bfseries 适用：}\par
      基础车桥耦合
    \end{PPTPanel}
  \end{column}
  \begin{column}{.245\textwidth}
    \begin{PPTPanel}{多刚体车辆}
      \centering
      \begin{tikzpicture}[x=.82cm,y=.82cm,>=Stealth]
        \fill[PanelFill] (.5,3.1) rounded rectangle (6.6,4.0);
        \draw[AuxiliaryC,line width=1.1pt] (.5,3.1) rounded rectangle (6.6,4.0);
        \node[font=\swjtufont{diagram}] at (3.55,3.55) {车体};
        \draw[IntersperseC,line width=1.1pt] (1.4,3.1)--(1.4,2.7);
        \draw[IntersperseC,line width=1.1pt] (5.6,3.1)--(5.6,2.7);
        \fill[AuxiliaryC!10] (1.0,2.2) rectangle (2.2,2.7);
        \fill[AuxiliaryC!10] (4.9,2.2) rectangle (6.1,2.7);
        \draw[AuxiliaryC] (1.0,2.2) rectangle (2.2,2.7);
        \draw[AuxiliaryC] (4.9,2.2) rectangle (6.1,2.7);
        \foreach \x in {1.2,2.0,5.1,5.9}{
          \filldraw[fill=white,draw=PrimaryC] (\x,1.7) circle (.25);
        }
        \draw[NormalTextC,line width=1.5pt] (.2,1.25)--(7.0,1.25);
      \end{tikzpicture}
      \PanelBulletsStart
        \item 多部件、多自由度
        \item 可描述空间耦合
      \PanelBulletsEnd
      \vfill
      {\swjtufont{body}\bfseries 适用：}\par
      复杂车辆与轨桥响应
    \end{PPTPanel}
  \end{column}
\end{columns}
\vfill
\begin{BottomBand}
  \centering
  模型选择：\textcolor{PrimaryC}{研究目标} \(\times\) \textcolor{AuxiliaryC}{响应精度}
  \(\times\) \textcolor{SecondaryAuxiliaryC}{计算效率}
\end{BottomBand}
\end{SlideBody}

% P4 — 简化车辆动力学模型的建立
\subsection{简单车辆模型的建立}

\begin{SlideBody}
\begin{columns}[T,onlytextwidth]
  \begin{column}{.36\textwidth}
    \begin{minipage}[t][10.65cm][s]{\linewidth}
      \begin{PPTPanel}{建模假设}
        \PanelBulletsStart
          \item 竖向二自由度
          \item 集中质量、刚体运动
          \item 线性悬挂与黏性阻尼
          \item 以静平衡位置为位移零点
        \PanelBulletsEnd
      \end{PPTPanel}
      \vfill
      \begin{PPTPanel}{主要参数}
        \PanelBulletsStart
          \item 簧上质量 \(m_s\)
          \item 簧下质量 \(m_u\)
          \item 悬挂刚度 \(k_s\)、阻尼 \(c_s\)
          \item 轮轨法向力增量 \(\Delta F_N\)
        \PanelBulletsEnd
      \end{PPTPanel}
    \end{minipage}
  \end{column}
  \begin{column}{.62\textwidth}
    \begin{PPTPanel}{二自由度竖向车辆模型}
      \centering
      \begin{tikzpicture}[x=1cm,y=1cm,>=Stealth,line cap=round,line join=round]
        \fill[PanelFill] (4.1,5.3) rounded rectangle (10.8,6.6);
        \draw[AuxiliaryC,line width=1.4pt] (4.1,5.3) rounded rectangle (10.8,6.6);
        \node[font=\swjtufont{larger}\bfseries,text=PrimaryC] at (7.45,5.95) {\(m_s\)};
        \draw[->,AuxiliaryC,line width=1.3pt] (11.3,6.2)--(11.3,5.35)
          node[right,font=\swjtufont{diagram},text=AuxiliaryC] {\(z_s\)};
        % spring
        \draw[IntersperseC,line width=1.6pt]
          (5.7,5.3)--(5.7,5.05)--++(.22,-.15)--++(-.44,-.22)
          --++(.44,-.22)--++(-.44,-.22)--++(.44,-.22)
          --++(-.22,-.15)--(5.7,3.8);
        % damper
        \draw[SecondaryAuxiliaryC,line width=1.4pt] (8.5,5.3)--(8.5,4.95);
        \draw[SecondaryAuxiliaryC,line width=1.4pt] (8.15,4.95) rectangle (8.85,4.45);
        \draw[SecondaryAuxiliaryC,line width=1.4pt] (8.5,4.45)--(8.5,3.8);
        \node[font=\swjtufont{diagram},text=IntersperseC] at (4.6,4.7) {\(k_s\)};
        \node[font=\swjtufont{diagram},text=SecondaryAuxiliaryC] at (9.2,4.7) {\(c_s\)};
        \fill[PanelFill] (4.1,2.5) rounded rectangle (10.8,3.8);
        \draw[AuxiliaryC,line width=1.4pt] (4.1,2.5) rounded rectangle (10.8,3.8);
        \node[font=\swjtufont{larger}\bfseries,text=PrimaryC] at (7.45,3.15) {\(m_u\)};
        \draw[->,AuxiliaryC,line width=1.3pt] (11.3,3.4)--(11.3,2.55)
          node[right,font=\swjtufont{diagram},text=AuxiliaryC] {\(z_u\)};
        % wheel and rail
        \filldraw[fill=white,draw=PrimaryC,line width=1.4pt] (7.45,1.65) circle (.55);
        \filldraw[fill=PanelFill,draw=AuxiliaryC,line width=1pt] (7.45,1.65) circle (.2);
        \draw[NormalTextC,line width=2pt] (3.3,.85)--(11.7,.85);
        \foreach \x in {3.7,4.7,...,11.0}{
          \draw[MutedTextC,line width=.8pt] (\x,.68)--++(.35,.34);
        }
        \draw[->,IntersperseC,line width=1.8pt] (7.45,.98)--(7.45,2.4);
        \node[font=\swjtufont{diagram}\bfseries,text=IntersperseC,anchor=west] at (8.0,1.45) {\(\Delta F_N\)};
        \node[font=\swjtufont{diagram},text=NormalTextC] at (7.4,.25) {轨道侧输入};
      \end{tikzpicture}
    \end{PPTPanel}
  \end{column}
\end{columns}
\vfill
\begin{BottomBand}
  \centering
  广义坐标：\(\displaystyle \mathbf q_v=\begin{bmatrix}z_s&z_u\end{bmatrix}^{\mathsf T}\)
  \qquad
  外力输入：\(\displaystyle \mathbf F_v=\begin{bmatrix}0&-\Delta F_N\end{bmatrix}^{\mathsf T}\)
\end{BottomBand}
\end{SlideBody}

% P5 — 车辆运动方程推导与矩阵表达
\subsection{车辆运动方程的建立}

\begin{SlideBody}
\begin{columns}[T,onlytextwidth]
  \begin{column}{.29\textwidth}
    \begin{PPTPanel}{受力分析}
      \centering
      \begin{tikzpicture}[x=.75cm,y=.75cm,>=Stealth,line cap=round,line join=round]
        \fill[PanelFill] (1.0,5.2) rectangle (7.3,6.3);
        \draw[AuxiliaryC,line width=1.1pt] (1.0,5.2) rectangle (7.3,6.3);
        \node[font=\swjtufont{diagram}] at (4.15,5.75) {\(m_s\)};
        \draw[IntersperseC,line width=1.4pt]
          (2.6,5.2)--(2.6,4.95)--++(.18,-.14)--++(-.36,-.18)
          --++(.36,-.18)--++(-.36,-.18)--(2.6,4.0);
        \draw[SecondaryAuxiliaryC,line width=1.4pt] (5.1,5.2)--(5.1,4.8);
        \draw[SecondaryAuxiliaryC,line width=1.4pt] (4.8,4.8) rectangle (5.4,4.35);
        \draw[SecondaryAuxiliaryC,line width=1.4pt] (5.1,4.35)--(5.1,4.0);
        \fill[PanelFill] (1.0,2.6) rectangle (7.3,4.0);
        \draw[AuxiliaryC,line width=1.1pt] (1.0,2.6) rectangle (7.3,4.0);
        \node[font=\swjtufont{diagram}] at (4.15,3.3) {\(m_u\)};
        \draw[->,IntersperseC,line width=1.5pt] (4.15,1.0)--(4.15,2.45);
        \node[font=\swjtufont{diagram},text=IntersperseC,anchor=west] at (4.45,1.55) {\(-\Delta F_N\)};
        \node[font=\swjtufont{diagram},text=IntersperseC] at (1.9,4.55) {\(k_s\)};
        \node[font=\swjtufont{diagram},text=SecondaryAuxiliaryC] at (5.9,4.55) {\(c_s\)};
        \draw[->,AuxiliaryC,line width=1.1pt] (7.55,5.95)--(7.55,5.0)
          node[right,font=\swjtufont{diagram}] {\(z_s\)};
        \draw[->,AuxiliaryC,line width=1.1pt] (7.55,3.7)--(7.55,2.75)
          node[right,font=\swjtufont{diagram}] {\(z_u\)};
        \node[font=\swjtufont{caption},text=MutedTextC,align=center] at (4.15,.35)
          {向下为正；接触力增量向上为正};
      \end{tikzpicture}
    \end{PPTPanel}
  \end{column}
  \begin{column}{.69\textwidth}
    \begin{minipage}[t][8.55cm][s]{\linewidth}
      \begin{PPTPanel}{簧上质量运动方程}
        \[
        m_s\ddot z_s+c_s(\dot z_s-\dot z_u)
        +k_s(z_s-z_u)=0
        \]
      \end{PPTPanel}
      \vfill
      \begin{PPTPanel}{簧下质量运动方程}
        \[
        m_u\ddot z_u+c_s(\dot z_u-\dot z_s)
        +k_s(z_u-z_s)=-\Delta F_N(t)
        \]
      \end{PPTPanel}
      \vfill
      \begin{PPTPanel}{矩阵标准形式}
        \[
        \mathbf M_v\ddot{\mathbf q}_v+
        \mathbf C_v\dot{\mathbf q}_v+
        \mathbf K_v\mathbf q_v=\mathbf F_v(t)
        \]
      \end{PPTPanel}
    \end{minipage}
  \end{column}
\end{columns}
\vfill
\begin{PPTPanel}{系统矩阵与外力向量}
  \centering
  \begin{tabular}{c c c c c}
    \(\displaystyle\mathbf M_v=\begin{bmatrix}m_s&0\\0&m_u\end{bmatrix}\) &
    \(\displaystyle\mathbf C_v=c_s\begin{bmatrix}1&-1\\-1&1\end{bmatrix}\) &
    \(\displaystyle\mathbf K_v=k_s\begin{bmatrix}1&-1\\-1&1\end{bmatrix}\) &
    \(\displaystyle\mathbf q_v=\begin{bmatrix}z_s\\z_u\end{bmatrix}\) &
    \(\displaystyle\mathbf F_v=\begin{bmatrix}0\\-\Delta F_N\end{bmatrix}\)
  \end{tabular}
\end{PPTPanel}
\end{SlideBody}

% P6 — 车辆运动方程与车桥耦合的联系
\section{车辆运动与车桥耦合}
\subsection{车辆方程与轮轨接口}

\begin{SlideBody}
\begin{columns}[T,onlytextwidth]
  \begin{column}{.35\textwidth}
    \begin{minipage}[t][9.8cm][s]{\linewidth}
      \begin{PPTPanel}{耦合变量}
        \PanelBulletsStart
          \item 轮对位移 \(z_u(t)\)
          \item 钢轨位移 \(z_r(x,t)\)
          \item 轨道不平顺 \(r(x)\)
          \item 轮轨力增量 \(\Delta F_N(t)\)
        \PanelBulletsEnd
      \end{PPTPanel}
      \vfill
      \begin{PPTPanel}{双向动力传递}
        \PanelBulletsStart
          \item 轮轨力经轨道结构传至桥梁
          \item 桥梁变形改变钢轨运动
          \item 钢轨运动反过来改变轮轨力
        \PanelBulletsEnd
      \end{PPTPanel}
    \end{minipage}
  \end{column}
  \begin{column}{.63\textwidth}
    \begin{PPTPanel}{车辆—轮轨—轨道—桥梁}
      \centering
      \begin{tikzpicture}[x=.92cm,y=.92cm,>=Stealth,line cap=round,line join=round]
        % train body and suspension
        \fill[PanelFill] (1.0,5.5) rounded rectangle (16.8,6.8);
        \draw[AuxiliaryC,line width=1.2pt] (1.0,5.5) rounded rectangle (16.8,6.8);
        \node[font=\swjtufont{diagram}\bfseries,text=PrimaryC] at (8.9,6.15) {车辆子系统};
        \draw[IntersperseC,line width=1.2pt] (4.5,5.5)--(4.5,4.8)--++(.18,-.12)--++(-.36,-.15)--++(.36,-.15)--(4.5,4.1);
        \draw[SecondaryAuxiliaryC,line width=1.2pt] (5.7,5.5)--(5.7,5.0)--(5.35,5.0)--(5.35,4.5)--(6.05,4.5)--(6.05,5.0)--(5.7,5.0)--(5.7,4.1);
        \fill[PanelFill] (3.0,3.4) rectangle (7.3,4.1);
        \draw[AuxiliaryC,line width=1.2pt] (3.0,3.4) rectangle (7.3,4.1);
        \node[font=\swjtufont{diagram},text=PrimaryC] at (5.15,3.75) {轮对};
        \filldraw[fill=white,draw=PrimaryC,line width=1.1pt] (5.15,2.8) circle (.42);
        % rail
        \draw[NormalTextC,line width=2.2pt] (.5,2.15)--(17.2,2.15);
        \draw[AuxiliaryC,line width=1.2pt,dashed] (.5,2.0)--(17.2,2.0);
        \node[font=\swjtufont{diagram},text=AuxiliaryC,anchor=west] at (13.4,2.5) {钢轨位移 \(z_r(x,t)\)};
        \draw[IntersperseC,decorate,decoration={snake,amplitude=.8mm,segment length=4mm},line width=1.1pt]
          (.5,2.9)--(17.2,2.9);
        \node[font=\swjtufont{diagram},text=IntersperseC,anchor=west] at (12.5,3.2) {不平顺 \(r(x)\)};
        % bridge beam and deformation
        \fill[PanelFill] (1.2,.3) rectangle (16.5,.8);
        \draw[AuxiliaryC,line width=1.2pt] (1.2,.3) rectangle (16.5,.8);
        \draw[AuxiliaryC,line width=1.2pt] (1.2,.8) .. controls (5.0,.3) and (12.0,.3) .. (16.5,.8);
        \draw[->,IntersperseC,line width=1.5pt] (5.15,3.2)--(5.15,2.05);
        \node[font=\swjtufont{diagram},text=IntersperseC,anchor=west] at (5.55,2.55) {\(\Delta F_N\)};
        \draw[->,AuxiliaryC,line width=1.3pt] (9.0,.95)--(9.0,1.75);
        \node[font=\swjtufont{diagram},text=AuxiliaryC,anchor=west] at (9.35,1.2) {桥梁变形 \(w(x,t)\)};
        \node[font=\swjtufont{diagram}\bfseries,text=PrimaryC] at (8.8,.05) {桥梁子系统};
      \end{tikzpicture}
    \end{PPTPanel}
  \end{column}
\end{columns}
\vfill
\begin{BottomBand}
  \centering
  \(\displaystyle \delta(t)=\delta_0+z_u(t)-z_r(x_c,t)-r(x_c),\quad x_c=vt\)
  \qquad
  \(\displaystyle F_N(t)=k_H[\delta(t)]_+^{3/2},\quad \Delta F_N=F_N-F_{N0}\)
\end{BottomBand}
\end{SlideBody}

% ============================================================
% Section 2: 最新研究进展
% ============================================================
\section{最新研究进展}
\subsection{考虑车辆速度效应的车桥耦合建模}

% P7 — Shi et al. (2026)
\begin{SlideBody}
\begin{columns}[T,onlytextwidth]
  \begin{column}{.34\textwidth}
    \begin{minipage}[t][10.65cm][s]{\linewidth}
      \begin{PPTPanel}{研究问题}
        传统车桥相互作用单元遗漏部分速度相关项。
      \end{PPTPanel}
      \vfill
      \begin{PPTPanel}{创新方法}
        \PanelBulletsStart
          \item 以移动接触点描述相对运动
          \item 纳入桥面斜率与曲率引起的速度效应
        \PanelBulletsEnd
      \end{PPTPanel}
      \vfill
      \begin{PPTPanel}{主要结论与拓展}
        \PanelBulletsStart
          \item 系统频率极值变化略有增强
          \item 瞬时阻尼比可能出现异常
          \item 重载高速工况需重点考虑
        \PanelBulletsEnd
      \end{PPTPanel}
    \end{minipage}
  \end{column}
  \begin{column}{.64\textwidth}
    \begin{minipage}[t][10.65cm][s]{\linewidth}
      \begin{PPTPanel}{移动接触点与桥梁几何效应}
        \centering
        \begin{tikzpicture}[x=.92cm,y=.92cm,>=Stealth,line cap=round,line join=round]
          \fill[PanelFill] (.4,4.3) rounded rectangle (17.6,5.6);
          \draw[AuxiliaryC,line width=1pt] (.4,4.3) rounded rectangle (17.6,5.6);
          \foreach \x in {1.0,2.5,4.0,5.5,7.0,8.5,10.0,11.5,13.0,14.5,16.0}{
            \draw[AuxiliaryC,line width=.7pt] (\x,4.65) rectangle ++(1.0,.48);
          }
          \node[font=\swjtufont{diagram}\bfseries,text=PrimaryC] at (9.0,5.2) {移动车辆};
          \draw[NormalTextC,line width=2pt] (.5,4.15)--(17.5,4.15);
          \draw[AuxiliaryC,line width=1.2pt] (.5,3.2) .. controls (5.0,2.7) and (12.5,2.7) .. (17.5,3.2);
          \draw[AuxiliaryC,dashed,line width=1pt] (.5,3.2)--(17.5,3.2);
          \draw[->,IntersperseC,line width=1.4pt] (8.7,6.0)--(11.0,6.0)
            node[midway,above,font=\swjtufont{diagram},text=IntersperseC] {\(v\)};
          \filldraw[fill=white,draw=IntersperseC,line width=1.3pt] (8.7,4.15) circle (.12);
          \draw[->,IntersperseC,line width=1.1pt] (8.7,4.05)--(8.7,3.45);
          \node[font=\swjtufont{diagram},text=IntersperseC,anchor=west] at (9.0,3.65) {\(x_c=vt\)};
          \node[font=\swjtufont{diagram},text=AuxiliaryC,anchor=west] at (11.0,2.75) {\(w(x,t)\)};
          \node[font=\swjtufont{diagram},text=PrimaryC,anchor=west] at (.8,2.45)
            {\(\theta=\partial w/\partial x,\quad \kappa=\partial^2w/\partial x^2\)};
        \end{tikzpicture}
      \end{PPTPanel}
      \vfill
      \begin{columns}[T,onlytextwidth]
        \begin{column}{.49\textwidth}
          \begin{PPTPanel}{时变频率}
            \centering
            \begin{tikzpicture}[x=.9cm,y=.75cm,>=Stealth]
              \draw[->,NormalTextC] (.2,.3)--(7.4,.3) node[right,font=\swjtufont{diagram}] {\(t\)};
              \draw[->,NormalTextC] (.2,.3)--(.2,2.5) node[above,font=\swjtufont{diagram}] {\(f(t)\)};
              \draw[AuxiliaryC,line width=1.5pt] plot[smooth] coordinates
                {(.4,.55)(1.2,.8)(2.0,1.9)(2.8,1.1)(3.8,.6)(4.8,1.0)(5.7,1.8)(6.8,.65)};
              \draw[MutedTextC,dashed,line width=1.2pt] plot[smooth] coordinates
                {(.4,.5)(1.2,.7)(2.0,1.35)(2.8,.95)(3.8,.55)(4.8,.8)(5.7,1.3)(6.8,.55)};
            \end{tikzpicture}
          \end{PPTPanel}
        \end{column}
        \begin{column}{.49\textwidth}
          \begin{PPTPanel}{瞬时阻尼比}
            \centering
            \begin{tikzpicture}[x=.9cm,y=.75cm,>=Stealth]
              \draw[->,NormalTextC] (.2,.3)--(7.4,.3) node[right,font=\swjtufont{diagram}] {\(t\)};
              \draw[->,NormalTextC] (.2,.3)--(.2,2.5) node[above,font=\swjtufont{diagram}] {\(\zeta(t)\)};
              \draw[AuxiliaryC,line width=1.5pt] plot[smooth] coordinates
                {(.4,1.0)(1.2,.7)(2.0,1.2)(2.8,.85)(3.8,1.75)(4.8,.9)(5.7,1.25)(6.8,.75)};
              \draw[MutedTextC,dashed,line width=1.2pt] plot[smooth] coordinates
                {(.4,.9)(1.2,.65)(2.0,1.0)(2.8,.8)(3.8,1.25)(4.8,.8)(5.7,1.0)(6.8,.7)};
            \end{tikzpicture}
          \end{PPTPanel}
        \end{column}
      \end{columns}
    \end{minipage}
  \end{column}
\end{columns}
\vfill
\begin{BottomBand}
  \centering
  \textcolor{PrimaryC}{研究启示：}移动接触运动学 \(\longrightarrow\) 速度相关项
  \(\longrightarrow\) 时变频率与阻尼
  \par
  {\swjtufont{caption}\color{MutedTextC}
  Shi et al., \emph{Mechanical Systems and Signal Processing}, 257 (2026), 114595.
  \href{https://doi.org/10.1016/j.ymssp.2026.114595}{DOI: 10.1016/j.ymssp.2026.114595}.
  曲线为概念示意，非论文原始数据。}
\end{BottomBand}
\end{SlideBody}

\subsection{三维车轨桥分区耦合与随机响应}

% P8 — Yuan & Luo (2026)
\begin{SlideBody}
\begin{columns}[T,onlytextwidth]
  \begin{column}{.34\textwidth}
    \begin{minipage}[t][9.45cm][s]{\linewidth}
      \begin{PPTPanel}{研究问题}
        \PanelBulletsStart
          \item 三维车轨桥系统时变耦合
          \item 随机响应呈非平稳、非高斯特征
        \PanelBulletsEnd
      \end{PPTPanel}
      \vfill
      \begin{PPTPanel}{创新方法}
        \PanelBulletsStart
          \item 车辆、轨道、桥梁非迭代分区
          \item 降维概率密度演化方法（DR-PDEE）
        \PanelBulletsEnd
      \end{PPTPanel}
      \vfill
      \begin{PPTPanel}{主要结论}
        \PanelBulletsStart
          \item 避免反复组装整体时变矩阵
          \item 以 10,000 组蒙特卡洛结果验证
          \item 可描述选定位移响应的概率演化
        \PanelBulletsEnd
      \end{PPTPanel}
    \end{minipage}
  \end{column}
  \begin{column}{.64\textwidth}
    \begin{minipage}[t][9.45cm][s]{\linewidth}
      \begin{PPTPanel}{三维车辆—轨道—桥梁分区模型}
        \centering
        \begin{tikzpicture}[x=.9cm,y=.9cm,>=Stealth,line cap=round,line join=round]
          % vehicle
          \fill[PanelFill] (.5,5.3) rounded rectangle (15.8,6.2);
          \draw[AuxiliaryC,line width=1.1pt] (.5,5.3) rounded rectangle (15.8,6.2);
          \node[font=\swjtufont{diagram}\bfseries,text=PrimaryC] at (8.1,5.75) {车辆子系统：7 个刚体};
          \foreach \x in {2.0,4.2,6.4,8.6,10.8,13.0,15.0}{
            \filldraw[fill=white,draw=PrimaryC,line width=.8pt] (\x,4.95) circle (.22);
          }
          % track
          \fill[PanelFill] (.5,3.65) rectangle (15.8,4.25);
          \draw[AuxiliaryC,line width=1.1pt] (.5,3.65) rectangle (15.8,4.25);
          \node[font=\swjtufont{diagram}\bfseries,text=PrimaryC] at (8.1,3.95) {轨道子系统：钢轨与轨道支承};
          % bridge
          \fill[AuxiliaryC!8] (.5,1.8) rectangle (15.8,2.5);
          \draw[AuxiliaryC,line width=1.1pt] (.5,1.8) rectangle (15.8,2.5);
          \node[font=\swjtufont{diagram}\bfseries,text=PrimaryC] at (8.1,2.15) {桥梁子系统：结构动力模型};
          % interfaces
          \draw[<->,IntersperseC,line width=1.6pt] (8.1,4.85)--(8.1,4.32);
          \node[font=\swjtufont{caption},text=IntersperseC,anchor=west] at (8.4,4.55) {时变轮轨接口};
          \draw[<->,SecondaryAuxiliaryC,line width=1.6pt] (8.1,3.6)--(8.1,2.55);
          \node[font=\swjtufont{caption},text=SecondaryAuxiliaryC,anchor=west] at (8.4,3.0) {轨桥接口};
          \node[font=\swjtufont{diagram},text=NormalTextC,anchor=west] at (.7,1.15)
            {车辆模型：51 自由度（含附加轮轨接口自由度）};
        \end{tikzpicture}
      \end{PPTPanel}
      \vfill
      \begin{PPTPanel}{随机响应的概率密度演化（方法示意）}
        \centering
        \begin{tikzpicture}[x=1cm,y=.8cm,>=Stealth]
          \draw[->,NormalTextC] (.1,.25)--(16.0,.25) node[right,font=\swjtufont{diagram}] {\(t\)};
          \foreach \x/\h/\col in {1.8/1.1/AuxiliaryC,5.4/1.55/SecondaryAuxiliaryC,9.2/1.25/IntersperseC,13.0/1.8/AlertedTextC}{
            \draw[->,NormalTextC] (\x-.7,.25)--(\x-.7,2.35);
            \draw[\col,line width=1.4pt]
              plot[domain=-1.35:1.35,samples=45]
              ({\x+\x/10*\pgfmathprintnumber{0}}, {.25+\h*exp(-(\x)^2)});
          }
          % Explicit bell curves, plotted in local coordinates for robustness.
          \draw[AuxiliaryC,line width=1.4pt] plot[domain=0:3.6,samples=55] (\x,{.25+1.15*exp(-((\x-1.8)^2)/.38)});
          \draw[SecondaryAuxiliaryC,line width=1.4pt] plot[domain=3.6:7.2,samples=55] (\x,{.25+1.45*exp(-((\x-5.4)^2)/.55)});
          \draw[IntersperseC,line width=1.4pt] plot[domain=7.2:10.8,samples=55] (\x,{.25+1.20*exp(-((\x-9.0)^2)/.48)});
          \draw[AlertedTextC,line width=1.4pt] plot[domain=10.8:14.8,samples=55] (\x,{.25+1.65*exp(-((\x-12.8)^2)/.60)});
          \node[font=\swjtufont{diagram},text=NormalTextC] at (8.0,2.55) {\(p(w,t)\)};
        \end{tikzpicture}
      \end{PPTPanel}
    \end{minipage}
  \end{column}
\end{columns}
\vfill
\begin{BottomBand}
  \centering
  车辆运动方程 \(\longrightarrow\) 子系统分区耦合 \(\longrightarrow\) 随机响应概率演化
  \par
  {\swjtufont{caption}\color{MutedTextC}
  Yuan \& Luo, \emph{Mechanical Systems and Signal Processing}, 259 (2026), 114850.
  \href{https://doi.org/10.1016/j.ymssp.2026.114850}{DOI: 10.1016/j.ymssp.2026.114850}.
  概率曲线为方法示意，非论文原始数据。}
\end{BottomBand}
\end{SlideBody}

% ============================================================
% Section 3: 个人问题及思考
% ============================================================
\section{个人问题及思考}
\subsection{速度相关项与随机耦合响应}

% P9 — 个人研究问题
\begin{SlideBody}
\begin{columns}[T,onlytextwidth]
  \begin{column}{.36\textwidth}
    \begin{minipage}[t][10.55cm][s]{\linewidth}
      \begin{PPTPanel}{研究问题}
        车辆速度相关项是否会改变车桥耦合随机响应的概率分布？
      \end{PPTPanel}
      \vfill
      \begin{PPTPanel}{研究方案}
        \PanelBulletsStart
          \item 基准模型：常规运动学表达
          \item 改进模型：显式保留速度相关项
          \item 保持随机输入与结构参数一致
        \PanelBulletsEnd
      \end{PPTPanel}
      \vfill
      \begin{PPTPanel}{评价指标}
        桥梁位移与加速度、轮轨力、响应分位数、计算耗时。
      \end{PPTPanel}
    \end{minipage}
  \end{column}
  \begin{column}{.62\textwidth}
    \begin{PPTPanel}{对照计算与敏感性分析}
      \centering
      \begin{tikzpicture}[x=.95cm,y=.95cm,>=Stealth,line cap=round,line join=round]
        % left model
        \fill[PanelFill] (.3,5.0) rounded rectangle (7.9,6.2);
        \draw[AuxiliaryC,line width=1.1pt] (.3,5.0) rounded rectangle (7.9,6.2);
        \node[font=\swjtufont{card-title}\bfseries,text=PrimaryC] at (4.1,5.8) {基准模型};
        \node[font=\swjtufont{diagram},text=NormalTextC] at (4.1,5.35) {常规接触运动学};
        % right model
        \fill[SoftGreen] (9.5,5.0) rounded rectangle (17.1,6.2);
        \draw[SecondaryAuxiliaryC,line width=1.1pt] (9.5,5.0) rounded rectangle (17.1,6.2);
        \node[font=\swjtufont{card-title}\bfseries,text=PrimaryC] at (13.3,5.8) {改进模型};
        \node[font=\swjtufont{diagram},text=NormalTextC] at (13.3,5.35) {显式保留速度项};
        \draw[->,AuxiliaryC,line width=1.5pt] (8.0,5.6)--(9.4,5.6);
        % shared inputs
        \fill[PanelFill] (2.0,3.5) rounded rectangle (15.5,4.35);
        \draw[AuxiliaryC!60!white] (2.0,3.5) rounded rectangle (15.5,4.35);
        \node[font=\swjtufont{diagram},text=PrimaryC] at (8.75,3.92)
          {同一车速、桥跨、车辆参数与轨道不平顺样本};
        \draw[->,AuxiliaryC,line width=1.3pt] (4.1,5.0)--(4.1,4.4);
        \draw[->,AuxiliaryC,line width=1.3pt] (13.3,5.0)--(13.3,4.4);
        % outputs
        \draw[->,IntersperseC,line width=1.3pt] (5.0,3.45)--(5.0,2.8);
        \draw[->,IntersperseC,line width=1.3pt] (12.5,3.45)--(12.5,2.8);
        \node[font=\swjtufont{diagram},text=IntersperseC,align=center] at (5.0,2.25)
          {桥梁位移与加速度\\轮轨力};
        \node[font=\swjtufont{diagram},text=IntersperseC,align=center] at (12.5,2.25)
          {概率分布与高分位数\\计算耗时};
        % bottom compare arrow
        \draw[<->,PrimaryC,line width=1.5pt] (5.9,1.1)--(11.6,1.1);
        \node[font=\swjtufont{diagram}\bfseries,text=PrimaryC] at (8.75,.55)
          {误差敏感性与适用边界};
      \end{tikzpicture}
    \end{PPTPanel}
  \end{column}
\end{columns}
\vfill
\begin{BottomBand}
  \centering
  \textcolor{PrimaryC}{个人思考：}在保证关键响应精度的前提下，识别速度效应不可忽略的工况范围。
\end{BottomBand}
\end{SlideBody}

\end{document}
